%% kapa: exam=1 round=29 year=2018 period=1 form=A subject=real_estate subject_ko=부동산학원론 seq=6 number=86
%% source: Q-Net 감정평가사 기출문제 (cst00309) artlSeq=5202808 file=extracted/1008/5202808/2018년 제29회 감정평가사 1차시험 기출문제/2018년 제29회 감정평가사 1차시험 기출문제/1교시 A형.pdf
%% conversion: kordoc 4.24.51 -> tools/gen1.py (auto); review=proofread-sonnet; flags=none
\begin{question}{86}
부동산시장의 효율성에 관한 설명으로 옳지 않은 것은? (단, 다른 조건은 고려하지 않음)
\begin{choices}
  \item 약성 효율적 시장은 현재의 시장가치가 과거의 추세를 충분히 반영하고 있는 시장이다.
  \item 준강성 효율적 시장은 어떤 새로운 정보가 공표되는 즉시 시장가치에 반영되는 시장이다.
  \item 강성 효율적 시장은 공표된 것이건 공표되지 않은 것이건 어떠한 정보도 이미 시장가치에 반영되어 있는 시장이다.
  \item 부동산시장은 주식시장이나 일반상품시장보다 더 불완전하고 비효율적이므로 할당 효율적일 수 없다.
  \item 부동산시장의 제약조건을 극복하는 데 소요되는 거래비용이 타 시장보다 부동산시장을 더 비효율적이게 하는 주요한 요인이다.
\end{choices}
\end{question}
\answer{4}
